\documentclass[11pt,a4paper]{article}
\usepackage{amsmath,amssymb,amsthm}
\usepackage{hyperref}
\usepackage{geometry}
\usepackage{booktabs}
\geometry{margin=1in}

\title{A Negative Definiteness Conjecture for a Weighted Cauchy Matrix\\
Arising from the Weil Explicit Formula}

\author{J.T.\ Vaughan}
\date{April 2026}

\newtheorem{theorem}{Theorem}
\newtheorem{conjecture}{Conjecture}
\newtheorem{proposition}{Proposition}
\newtheorem{lemma}{Lemma}
\newtheorem{definition}{Definition}
\newtheorem{remark}{Remark}
\newtheorem{observation}{Observation}

\DeclareMathOperator{\diag}{diag}
\DeclareMathOperator{\tr}{tr}

\begin{document}

\maketitle

\begin{abstract}
We present an open conjecture about the negative definiteness of a parametric family of real symmetric matrices with explicit Cauchy--Loewner structure. The matrix $M_{\mathrm{odd}}(\lambda)$ has dimension $N \times N$ with entries
\[
M_{jk} = \frac{2(jB_k - kB_j)}{j^2 - k^2}, \quad j \neq k, \qquad M_{jj} = a_j + \frac{B_j}{j},
\]
where the sequences $a_n$ and $B_n$ are determined by the Weil explicit formula for the Riemann zeta function. The conjecture---that $M_{\mathrm{odd}}(\lambda) < 0$ for all $\lambda$---is equivalent to the Riemann Hypothesis through the Connes--Weil positivity framework. We provide the precise definitions, closed-form entry formulas, extensive computational evidence, and structural properties (knife-edge sensitivity, arithmetic specificity) that constrain any proof strategy. The purpose of this note is to formulate the conjecture as a pure matrix analysis problem, accessible to specialists in Cauchy matrix theory without requiring background in analytic number theory.
\end{abstract}

%======================================================================
\section{The Matrix}
%======================================================================

Fix a parameter $\lambda > 1$ and set $L = \log(\lambda^2)$, $N = \max(15, \lceil 6L \rceil)$. Define two sequences for $n = 1, 2, \ldots, N$:

\begin{definition}[Diagonal sequence $a_n$]\label{def:an}
\begin{equation}\label{eq:an}
a_n = \underbrace{\mathrm{wr}(n, L)}_{\text{archimedean}} + \underbrace{\sum_{\substack{p^k \leq \lambda^2 \\ p \text{ prime}}} \frac{\log p}{p^{k/2}} \cdot \frac{2(L - k\log p)}{L} \cdot \cos\!\Bigl(\frac{2\pi n k \log p}{L}\Bigr)}_{\text{prime sum}},
\end{equation}
where the archimedean function $\mathrm{wr}(n, L)$ is defined by
\begin{equation}\label{eq:wr}
\mathrm{wr}(n, L) = \gamma + \log\!\Bigl(\frac{4\pi(e^L - 1)}{e^L + 1}\Bigr) + \int_0^\infty \frac{e^{x/2}\,\omega_n(x,L) - 2}{e^x - e^{-x}}\,dx,
\end{equation}
with $\omega_n(x, L) = 2\cos(2\pi n x/L)\cdot\mathbf{1}_{[0,L]}(x)$, $\gamma$ the Euler--Mascheroni constant, and asymptotic behavior $\mathrm{wr}(n,L) = C(L) - \log n + O(1/n^2)$ where $C(L) = \gamma + \log(4\pi(e^L-1)/(e^L+1)) > 0$.
\end{definition}

\begin{definition}[Off-diagonal sequence $B_n$]\label{def:Bn}
\begin{equation}\label{eq:Bn}
B_n = \underbrace{\alpha(n, L)}_{\text{archimedean}} + \underbrace{\sum_{\substack{p^k \leq \lambda^2 \\ p \text{ prime}}} \frac{\log p}{\pi\, p^{k/2}} \cdot \sin\!\Bigl(\frac{2\pi n k \log p}{L}\Bigr)}_{\text{prime sum}},
\end{equation}
where the archimedean function $\alpha(n, L)$ is defined by
\begin{equation}\label{eq:alpha}
\alpha(n, L) = \frac{1}{\pi}\biggl[\frac{e^{-L/2}\operatorname{Im}\bigl(\frac{2L}{L + 4\pi i n}\,{}_2F_1(1, a, a{+}1, e^{-2L})\bigr) + \frac{1}{2}\operatorname{Im}\,\psi(a)}{1}\biggr],
\end{equation}
with $a = \tfrac{1}{4} + \tfrac{i\pi n}{L}$, $\psi$ the digamma function, and ${}_2F_1$ the Gauss hypergeometric function. The sequence satisfies $\alpha(-n, L) = -\alpha(n, L)$ (odd symmetry) and $\alpha(n, L) = O(1/n)$.
\end{definition}

\begin{definition}[The matrix $M_{\mathrm{odd}}$]\label{def:Modd}
The real symmetric $N \times N$ matrix $M_{\mathrm{odd}}$ has entries:
\begin{align}
(M_{\mathrm{odd}})_{jj} &= a_j + \frac{B_j}{j}, \quad j = 1, \ldots, N, \label{eq:diag}\\
(M_{\mathrm{odd}})_{jk} &= \frac{2(j B_k - k B_j)}{j^2 - k^2}, \quad j \neq k. \label{eq:offdiag}
\end{align}
\end{definition}

\begin{remark}
The matrix $M_{\mathrm{odd}}$ is the restriction of a larger $(2N+1) \times (2N+1)$ Cauchy--Loewner matrix to its odd-parity subspace under $n \mapsto -n$. The full matrix has the simpler form $M_{nm} = a_n\, \delta_{nm} + (B_m - B_n)/(n - m)$, and the off-diagonal formula~\eqref{eq:offdiag} arises from the parity projection. In the notation of Pushnitski--Treil~\cite{PT2025}, this is a weighted Cauchy matrix with displacement rank~$2$.
\end{remark}

%======================================================================
\section{The Conjecture}
%======================================================================

\begin{conjecture}[Negative definiteness]\label{conj:main}
For all $\lambda \geq \sqrt{2}$ (equivalently $L \geq \log 2$) and all $N \geq N_0(\lambda)$, the matrix $M_{\mathrm{odd}}(\lambda, N)$ is negative definite.
\end{conjecture}

\begin{theorem}[Consequence]\label{thm:consequence}
Conjecture~\ref{conj:main}, combined with a verified barrier condition on the complementary (even-parity) block, implies the Weil positivity criterion $Q_W \geq 0$, which is equivalent to the Riemann Hypothesis. The proof of this implication appears in~\cite{Vaughan2026}.
\end{theorem}

%======================================================================
\section{Structural Properties}
%======================================================================

\subsection{The diagonal is logarithmic}

The archimedean component of $a_n$ satisfies $\mathrm{wr}(n, L) = C(L) - \log n + O(1/n^2)$ where $C(L) > 0$ for all~$L$. This logarithmic decay is the fingerprint of the Gamma function $\Gamma(s/2)$ in the functional equation of~$\zeta$, and corresponds to the trivial zeros of~$\zeta$ at $s = -2, -4, -6, \ldots$

\begin{observation}[Knife-edge]\label{obs:knife}
Scaling the archimedean diagonal $\mathrm{wr}(n,L)$ by a factor~$\alpha$ destroys negative definiteness for $\alpha \neq 1$:
\begin{center}
\begin{tabular}{rcc}
\toprule
$\alpha$ & \#positive eigenvalues & $\lambda_{\max}(M_{\mathrm{odd}})$ \\
\midrule
$0.95$ & $27$ & $+0.175$ \\
$\mathbf{1.00}$ & $\mathbf{0}$ & $\mathbf{-1.6 \times 10^{-7}}$ \\
$1.05$ & $6$ & $+0.160$ \\
\bottomrule
\end{tabular}
\end{center}
The $-\log n$ decay rate is the \emph{unique} rate that produces negative definiteness. Replacing $-\log n$ with $-n^{\alpha}$ for any $\alpha > 0$ fails.
\end{observation}

\subsection{Arithmetic specificity}

\begin{observation}[Cram\'er model]\label{obs:cramer}
Replacing the actual prime numbers with random integers drawn from the Cram\'er model (each $n$ is ``prime'' with probability $1/\log n$) produces a matrix that is \emph{never} negative definite: $0$ out of $200$ trials yield $M_{\mathrm{odd}} < 0$. The minimum $\lambda_{\max}(M_{\mathrm{odd}})$ across Cram\'er trials is $+4.26$, compared to the actual value of~$-1.6 \times 10^{-7}$---a gap of $26$ million. Similarly, $0/1000$ random diagonal perturbations and $0/1000$ log-shaped diagonal perturbations yield negative definiteness.
\end{observation}

\subsection{The critical direction}

\begin{observation}[Ten-figure cancellation]\label{obs:cancel}
The maximum eigenvalue of $M_{\mathrm{odd}}$ at $\lambda^2 = 1000$ is $-1.6 \times 10^{-7}$, achieved on an eigenvector $v \approx -0.54\, e_1 + 0.84\, e_2$ concentrated in the first two components. The Rayleigh quotient decomposes as:
\begin{center}
\begin{tabular}{lr}
\toprule
Component & Contribution \\
\midrule
Prime off-diagonal & $-1.534$ \\
Archimedean diagonal & $+1.565$ \\
Alpha correction & $-0.031$ \\
\midrule
Total & $-1.6 \times 10^{-7}$ \\
\bottomrule
\end{tabular}
\end{center}
\end{observation}

\subsection{Convergence with truncation}

The maximum eigenvalue of $M_{\mathrm{odd}}$ is $N$-converged: it stabilizes from $-1.77 \times 10^{-7}$ at $N = 15$ to $-1.56 \times 10^{-7}$ at $N = 150$.

\subsection{Sylvester leading minors}

All leading principal minors of $M_{\mathrm{odd}}$ satisfy the Sylvester criterion for negative definiteness. The $k$-th minor has magnitude decreasing as $\sim 10^{-5k}$, indicating extreme cancellation at every Schur step.

%======================================================================
\section{What a Proof Would Need}
%======================================================================

Every standard matrix-analytic bound fails by many orders of magnitude:

\begin{itemize}
\item \textbf{Gershgorin:} The Gershgorin disk for row~$1$ is $[-64, +40]$, failing to exclude positive eigenvalues by a factor of $10^8$.
\item \textbf{Diagonal dominance:} The off-diagonal row sums exceed the diagonal entries by factors of $4$--$37$.
\item \textbf{Trace + determinant:} These constrain but do not determine the signature.
\item \textbf{Loewner inertia (Bharali--Holtz):} Their theorem gives Lorentzian signature for $f(x) = x^r$ with $1 < r < 2$, but $B_n \sim 1/n$ is not in this function class.
\item \textbf{Cauchy interlacing:} Each prime contributes a full-rank perturbation, not rank-$1$, precluding eigenvalue tracking.
\end{itemize}

A proof likely requires exploiting the specific algebraic relationship between $a_n$ and $B_n$ imposed by the Weil explicit formula. Both sequences are determined by the same set of prime numbers, and the computational evidence (Section~\ref{sec:evidence} and Observation~\ref{obs:cramer}) strongly suggests that this coupling is essential for negative definiteness.

\begin{remark}[The functional equation constraint]
The functional equation $\xi(s) = \xi(1-s)$ forces $a_n$ to be even and $B_n$ to be odd in~$n$. More importantly, it implies that the test function $h_{\mathrm{pair}}(\delta, \gamma) = h(\delta + i\gamma) + h(-\delta + i\gamma)$ satisfies $\partial h_{\mathrm{pair}} / \partial \delta \big|_{\delta = 0} = 0$ for all~$\gamma$---the critical line $\mathrm{Re}(s) = 1/2$ is automatically a stationary point of the zero contribution. The second derivative is positive (minimum), meaning deviations from the critical line are destabilizing.
\end{remark}

%======================================================================
\section{Computational Evidence}\label{sec:evidence}
%======================================================================

Negative definiteness has been verified at the following parameter values (and hundreds of others):

\begin{center}
\begin{tabular}{rrrrr}
\toprule
$\lambda^2$ & $N$ & $\lambda_{\max}(M_{\mathrm{odd}})$ & $\tr(M_{\mathrm{odd}})$ & $|\lambda_{\max}|/|\tr|$ \\
\midrule
$50$ & $23$ & $-1.86 \times 10^{-7}$ & $-24.9$ & $7.5 \times 10^{-9}$ \\
$200$ & $32$ & $-1.74 \times 10^{-7}$ & $-40.2$ & $4.3 \times 10^{-9}$ \\
$1000$ & $41$ & $-1.61 \times 10^{-7}$ & $-63.7$ & $2.5 \times 10^{-9}$ \\
$5000$ & $51$ & $-1.48 \times 10^{-7}$ & $-112.0$ & $1.3 \times 10^{-9}$ \\
$50000$ & $65$ & $-1.27 \times 10^{-7}$ & $-273.1$ & $4.7 \times 10^{-10}$ \\
$200000$ & $73$ & $-1.14 \times 10^{-7}$ & $-437.2$ & $2.6 \times 10^{-10}$ \\
\bottomrule
\end{tabular}
\end{center}

The ratio $|\lambda_{\max}|/|\tr|$ decreases, showing that the trace becomes an increasingly negligent bound on the maximum eigenvalue. The maximum eigenvalue itself appears to scale as $\sim e^{-0.17 L}$, approaching zero but remaining negative at all tested parameters.

All Python scripts and numerical data are available in the project repository.

%======================================================================
\section{Connections to Existing Theory}
%======================================================================

\begin{itemize}
\item The full $(2N+1) \times (2N+1)$ matrix $M$ satisfies the displacement equation $[D, M] = \mathbf{1}\mathbf{B}^T - \mathbf{B}\mathbf{1}^T$ where $D = \diag(n)$, giving displacement rank~$2$ in the sense of Kailath--Lev-Ari~\cite{Kailath1979}. However, the odd-parity projection inflates the displacement rank to~$14$.

\item The matrix is connected to the spectral triples of Connes~\cite{CC2025}, who constructs self-adjoint operators from rank-$1$ perturbations of scaling operators. The matrix entries coincide with Connes' ``$\tau$-matrix'' from~\cite{CC2025} to $10^{-15}$ precision.

\item The Huh--Br\"and\'en theory of Lorentzian polynomials~\cite{BH2020} characterizes matrices with at most one positive eigenvalue. The full matrix~$M$ (before parity projection) has signature $(1, 2N)$---Lorentzian---while $M_{\mathrm{odd}}$ has signature $(0, N)$---negative definite. Both properties are required for the Riemann Hypothesis.
\end{itemize}

%======================================================================
\section*{Acknowledgments}
%======================================================================

All numerical results are reproducible from the Python scripts in the project repository at \texttt{github.com/vaughanjt/RiemannHypothesisProof}. Negative definiteness has been rigorously verified at $134$ parameter values using interval arithmetic (python-flint, $256$-bit precision) with the Sylvester criterion.

%======================================================================
\begin{thebibliography}{99}
%======================================================================

\bibitem{BH2020}
P.~Br\"and\'en and J.~Huh,
\emph{Lorentzian polynomials},
Ann.\ of Math.\ \textbf{192} (2020), 821--891.

\bibitem{CC2025}
A.~Connes and C.~Consani,
\emph{Zeta spectral triples},
preprint (2025), arXiv:2511.22755.

\bibitem{Kailath1979}
T.~Kailath, S.~Kung, and M.~Morf,
\emph{Displacement ranks of matrices and linear equations},
J.\ Math.\ Anal.\ Appl.\ \textbf{68} (1979), 395--407.

\bibitem{PT2025}
A.~Pushnitski and S.~Treil,
\emph{The spectral map for weighted Cauchy matrices is an involution},
preprint (2025), arXiv:2504.18707.

\bibitem{Vaughan2026}
J.T.~Vaughan,
\emph{The Lorentzian Weil Matrix Conjecture: A Reduction of the Riemann Hypothesis to Cauchy--Loewner Matrix Theory},
preprint (2026).

\end{thebibliography}

\end{document}
